\section{推荐系统 vs 生成系统：没有中间商赚差价}

本章进入全片最核心的平台理论。推荐系统的本质是库存分配：平台把已有内容和用户偏好匹配起来，掌握排序、分发和流量定价。生成系统的本质是按需生产：用户的需求直接进入生产端，生产端生成内容再直接给用户，分配环节被内化到生产系统里。

\begin{figure}[H]
\centering
\includegraphics[width=0.94\textwidth]{recommendation-vs-generation.png}
\caption{推荐系统 vs 生成系统：库存分配与按需生产的根本差异。自制概念图，依据 02:38:11--02:50:34 对谈内容整理。}
\end{figure}

\begin{knowledgebox}{读图：中间商是“库存分配权”}
推荐系统从已有库存中匹配内容，平台掌握分发权；生成系统按用户需求即时生产，分配内化到生产。所谓“没有中间商赚差价”，不是没有平台，而是旧平台的分发权不再是唯一价值兑现方式。
\end{knowledgebox}

\subsection{推荐系统为什么是中间商}

本节先解释“中间商”这个词的机制含义。它不是道德批评，而是说平台掌握库存分配、排序和商业化权力，因此能在生产者和消费者之间抽取价值。

互联网平台的权力来自双边关系：创作者生产内容，消费者消费内容，平台负责匹配并抽取价值。搜索引擎、推荐引擎、电商平台都在不同程度上做中间商。创作者投入成本生产内容，但同样内容由不同账号发布，流量可能完全不同，这说明分配机制不只取决于内容本身，也取决于平台权重、账号历史、算法目标和商业规则。

\begin{warningbox}{不要把“平台是中间商”理解为平台没有价值}
平台确实创造了发现、匹配、反作弊、支付、社区和规模化分发价值。问题是生成系统会削弱“库存分配”这一项的相对稀缺性，迫使平台转向生产系统、信任网络或其他价值层。
\end{warningbox}

\subsection{生成系统如何内化分配}

生成系统不是从库存里找一个最接近用户需求的内容，而是根据用户意图直接生产内容。它仍然有分配，但分配发生在生成过程里：用户的需求、偏好、context、recipe 和系统能力共同决定输出。这样一来，平台不再唯一控制内容是否被看见，用户意图成为更直接的生产输入。

\begin{figure}[H]
\centering
\includegraphics[width=0.96\textwidth]{power-shift-to-consumer.png}
\caption{权力向消费者端渗透：信息商品链条从平台分配走向用户意图驱动。自制概念图，依据 02:25:49--02:38:11 对谈内容整理。}
\end{figure}

\begin{knowledgebox}{读图：消费者不只是接收端}
传统链条中，生产者生产、平台分发、消费者接收。生成系统让消费者意图进入生产端，用户不只是被动消费内容，而是在消费时参与生产内容。
\end{knowledgebox}

\subsection{本章小结}

推荐系统和生成系统的根本差异，是库存分配和按需生产的差异。平台不会消失，但旧平台靠分发权抽取价值的方式会被挑战。

